% Generated draft: requires mathematical and source review.
\documentclass[11pt]{article}
\usepackage[margin=25mm]{geometry}
\usepackage{fontspec}
\setmainfont{Latin Modern Roman}
\usepackage{amsmath,amssymb,mathtools,unicode-math}
\setmathfont{Latin Modern Math}
\usepackage{enumitem,needspace,xurl,hyperref}
\setlength{\parindent}{0pt}
\setlength{\parskip}{5pt}
\setlength{\emergencystretch}{3em}
\newcommand{\N}{\mathbb{N}}
\newcommand{\Z}{\mathbb{Z}}
\newcommand{\Q}{\mathbb{Q}}
\newcommand{\R}{\mathbb{R}}
\newcommand{\C}{\mathbb{C}}
\newcommand{\F}{\mathbb{F}}
\newcommand{\A}{\mathbb{A}}
\newcommand{\PP}{\mathbb{P}}
\newcommand{\E}{\mathbb{E}}
\newcommand{\eps}{\varepsilon}
\newcommand{\dd}{\,\mathrm{d}}
\begin{document}
\section*{UnsolvedMath ID: 6400007}
\textbf{Source problem code:} catalogue entry
\subsection{Definitions and mathematical statement}
Let $M$ be a compact connected smooth seven-manifold, with the usual convention of no boundary. The compact exceptional group $G_2\subset SO(7)$ is the stabilizer of a positive three-form on $\mathbb R^7$. For a Riemannian metric $g$, its full holonomy group $\operatorname{Hol}(M,g)$ is generated by parallel transports around based loops.

Characterize those $M$ for which there exists a smooth Riemannian metric satisfying
\[
\operatorname{Hol}(M,g)=G_2
\]
up to conjugacy in $SO(7)$. Equality is required; containment in $G_2$ is a weaker condition.

A related structure formulation uses a positive three-form $\varphi$ inducing $g_\varphi$ and satisfying $d\varphi=0$ and $d(*_{\varphi}\varphi)=0$. Such a torsion-free $G_2$ structure gives holonomy contained in $G_2$; the proper-subgroup cases must still be excluded to address exact holonomy.

The supplied summary lists spin and $b_1(M)=0$ among necessary conditions, mentions unspecified constraint relations and many examples, and also claims ``2-connected for exact holonomy.'' Two-connected means $\pi_1(M)=\pi_2(M)=0$; this strong claim is neither an assumption of the question nor accompanied by a supporting theorem in the record. It is not inserted as a restriction on the manifolds to be classified. The target remains a complete characterization of all compact seven-manifolds admitting exact $G_2$ holonomy.
\subsection{Short English statement}
Which compact smooth seven-manifolds admit a metric whose full holonomy is exactly $G_2$? A torsion-free $G_2$ structure with smaller holonomy does not by itself answer this classification question.
\subsection{Sources}
\begin{itemize}
\item Simon Donaldson, \emph{Some Problems in Differential Geometry and Topology} (2008), explicit question/program 7 in the supplied author PDF, \url{https://wwwf.imperial.ac.uk/~skdona/NONLINEARITYPROBLEMS.PDF}; accessed 2026-07-29, PDF-text-normalized extraction.
\item Local registry \texttt{data/unsolvedmath/problems.json}, permanent UnsolvedMath ID $6400007$, source code catalogue entry, \emph{Topology of Compact Manifolds with Holonomy $G_2$}. Dataset attribution: Ulam AI and the UnsolvedMath Contributors.
\item The research summary gives topological-condition and example claims without individual references or explicit constraint formulas.
\end{itemize}
\end{document}
